\subsection{Limitations and Future Work}

The real-flight evidence has defined limits. UAV-SEAD labels describe runtime anomalies and fault domains, not faulty source modules. The fixed chronological reproduction permits a source log to span partitions; only the leave-log-out protocol enforces complete-log isolation. The channel dictionary was selected by unlabelled corpus coverage, whereas each protocol's scaler and supervised models use its training data. External validation is binary only, and Uncategorized screening lacks temporal labels.

The explanation and architecture layers are associative. The semantic mapping is author-defined and not independently expert-adjudicated; External and Global Position consistency is weak. The uORB parser's perfect local agreement is an internal shared-vocabulary check, not an independent gold audit. Conservative path exclusion does not reconstruct a target-specific compiled graph. The real-flight architecture analysis covers the 38.7\% matching-commit anomalous subset, while propagation depends on topic degree and publisher/subscriber semantics. Neither this representation nor its masking tests measure engineering inspection time or causal necessity.

The paired-replay extension requires an additional healthy reference from the same source flight and evaluates only four fixed mutation mechanisms at one firmware commit. Its three development and three transfer source flights provide limited independent variation; extra normal replays do not add new source flights. Although replay parameters were frozen before the first transfer evaluation, the historical P2 scaler had already encountered the transfer sources, and subsequent calibration variants used already inspected transfer flights for exploratory assessment. These results are not a fully untouched-pipeline generalization test. Mutation-effect checks failed for one Commander trial in each source set; those trials remain in the primary denominators. Missing or widely separated observations yield unavailable evidence rather than normality. Full-log interpolation and cadence checks, and whole-run evidence aggregation, also preclude an online or real-time claim.

Future work should first isolate healthy EKF2 replay divergence and improve observation-aware residuals without using mutation-specific prediction rules. Any revised method requires separate development and calibration flights followed by a fully untouched source-flight evaluation, including the scaler. Independent semantic adjudication, per-commit graphs, and runtime topic provenance would address the separate limitations of static module evidence. These are extensions of the present evidence framework rather than prerequisites for interpreting its recorded-data classification and explanation results.
